\section{Combining the existing spectral and overlap controls on the original source}
\label{sec:ACM}
This attachment records the exact comparison already established in the
coordinating source intake. The complete earlier AW proof is retained in
this reader; SP reproduces its spectra and final condition-number bound
and adds the stated overlap refinement. Neither bound is discarded.

Fix the actual full packet $h$, tensor order $k\ge1$, and its original
cyclic polynomial $\chi$ of degree $q\ge1$. The common coefficient space
is $H=\mathcal P_{2q}$. Its integration coordinate is $S=k/2+iy$.
Write $x\in[0,1]$ for metric interpolation and reserve $u$ for the
geometric deformation parameter of BC. The exact AW/SP coordinate map is
\[
 y=u_{\rm AW},\qquad x=t_{\rm AW},\qquad
 H_{\rm AW}\xrightarrow{\operatorname{id}}H_{\rm SP}.
 \tag{ACM1}
\]
This identity leaves every polynomial coefficient unchanged. The
arithmetic density is $w_h^{*k}$ in both calculations. For the reference
density, the original constant $c_a=2^{1-2a}\Gamma(2a)$ gives
$c_{1/4}=\sqrt{2\pi}$, so the exact convolution formula is
\[
 r_{1/4}^{*k}(y)=\frac{c_{1/4}^k}{c_{k/4}}r_{k/4}(y)
 =\frac{(\sqrt{2\pi})^k}{c_{k/4}}
       \frac{|\Gamma(k/4+iy/2)|^2}{2\pi}.
 \tag{ACM2}
\]
Consequently the densities, their masses and all monomial integrals
coincide. In particular $M_{\rm AW}(x)=M_{\rm SP}(x)=M_x$ on $H$,
and $M_{N,\rm AW}=I_N^*M_xI_N$ on the original subspace
$I_N:\mathcal P_N\hookrightarrow H$.

AW's relation map has codomain $\mathcal P_N$, whereas SP incorporates
its inclusion into $H$. Their literal equality is
\[
 B_N^{\rm SP}=I_NB_N^{\rm AW}.
 \tag{ACM3}
\]
Substitution into the orthogonal projector formula gives
\[
 \begin{aligned}
 P_N&=I_N(I_N^*M_xI_N)^{-1}I_N^*M_x,\\
 Q_N&=I_NB_N^{\rm AW}
       ((B_N^{\rm AW})^*I_N^*M_xI_NB_N^{\rm AW})^{-1}
       (B_N^{\rm AW})^*I_N^*M_x.
 \end{aligned}
 \tag{ACM4}
\]
The two definitions of each projector therefore agree, including
$Q_{q-1}=0$ for the zero relation domain. Put
\[
 \begin{aligned}
 U_x&=Q_{2q-1}+Q_{2q}-Q_q=\mathcal R_x^{\rm SP},\\
 W_x&=(P_{2q-1}-P_{q-1})+(P_{2q}-P_q)=\mathcal P_x^{\rm SP},\\
 C_x&=M_x^{-1}\dot M_x=T_x^{\rm SP}.
 \end{aligned}
 \tag{ACM5}
\]
These are AW's original $U,W,C$. Its independently constructed
source isometry $T^{\rm iso}_x$ satisfies
\[
 (T^{\rm iso}_x)^*M_xT^{\rm iso}_x=M_x,\qquad
 U_x=T^{\rm iso}_xW_x(T^{\rm iso}_x)^{-1}.
 \tag{ACM6}
\]
Thus cyclicity of the same finite trace proves
\[
 \operatorname{Tr}((U_x-W_x)C_x)
 =\operatorname{Tr}\bigl(W_x(T^{\rm iso}_x)^{-1}
                             [C_x,T^{\rm iso}_x]\bigr).
 \tag{ACM7}
\]
The original observation is still
$\mathcal O=\sigma_h^{\otimes k}\eta J_{2q}$. Its exact difference
and kernel after the isometry remain
\[
 \begin{aligned}
 \mathcal OT^{\rm iso}_x-\mathcal O
   &=\sigma_h^{\otimes k}\eta J_{2q}(T^{\rm iso}_x-I),\\
 \ker(\mathcal OT^{\rm iso}_x)
   &=(T^{\rm iso}_x)^{-1}(\chi\mathcal P_q).
 \end{aligned}
 \tag{ACM8}
\]
The first equality is expansion. The second follows because
$\sigma_h^{\otimes k}$ and $\eta$ are the original injections and
$\ker J_{2q}=\chi\mathcal P_q$. This retains the actual cohomology
defect of the isometry rather than assigning it an intertwining property.

Let $F(x)$ be the original logarithmic four-endpoint ratio and
$\Delta_{h,k}=F(1)-F(0)$. Its signed identity is
$F'(x)=\operatorname{Tr}((U_x-W_x)C_x)$ by (AW6), or equivalently
(SP4)--(SP5) and (ACM5). Both $U_x,W_x$ have the complete spectrum
$0$ with multiplicity $q$, $1$ with multiplicity $2$, and $2$ with
multiplicity $q-1$. All statements below also cover $q=1$.

Write $c_1(x)\le\cdots\le c_{2q+1}(x)$ for the eigenvalues of $C_x$
in the original $M_x$ metric and $d_x=c_{2q+1}(x)-c_1(x)$.
Let $s_x=\dim(\operatorname{ran}U_x\cap\operatorname{ran}W_x)$.
The two already proved pointwise bounds, on this identical source, are
\[
 \begin{aligned}
 S_{\rm AW}(x)&=c_{q+2}(x)-c_q(x)
       +2\sum_{j=1}^{q-1}(c_{2q+2-j}(x)-c_j(x)),\\
 S_{\rm SP}(x)&=\frac{d_x}{2}\min\left\{
       4q-2s_x,\sqrt{(2q+1)(8q-4-2\operatorname{Tr}(U_xW_x))}
                                    \right\}.
 \end{aligned}
 \tag{ACM9}
\]
For completeness, their simultaneous application follows directly as
follows. An orthogonal projection $P$ of rank $r$ satisfies
$\sum_{j=1}^r c_j\le\operatorname{Tr}(PC_x)
\le\sum_{j=2q+2-r}^{2q+1}c_j$: in an orthonormal eigenbasis of
$C_x$, its diagonal entries lie in $[0,1]$ and have sum $r$, so
moving their mass to the lowest or highest eigenvalues proves each
inequality. Each of $U_x,W_x$ is the sum of its range projection
of rank $q+1$ and its eigenvalue-two projection of rank $q-1$.
Apply these bounds to both sums and subtract their extrema to obtain
$|F'(x)|\le S_{\rm AW}(x)$, with every multiplicity retained.

The range dimensions give $s_x\ge1$. If $E_x$ projects onto the
intersection, then $U_x-E_x$ and $W_x-E_x$ are positive, each of
trace $2q-s_x$. Therefore
$\|U_x-W_x\|_1\le4q-2s_x$. The same complete spectra give
\[
 \operatorname{Tr}(U_x-W_x)=0,\qquad
 \operatorname{Tr}((U_x-W_x)^2)
       =8q-4-2\operatorname{Tr}(U_xW_x).
 \tag{ACM10}
\]
Cauchy--Schwarz on all $2q+1$ real eigenvalues gives the other
trace-norm bound in (ACM9). Subtract the midpoint
$(c_1+c_{2q+1})I/2$ from $C_x$ in the trace pairing; the first
identity in (ACM10) leaves $F'$ unchanged. Diagonalizing $U_x-W_x$
then bounds the pairing by $d_x\|U_x-W_x\|_1/2$. This proves
$|F'(x)|\le S_{\rm SP}(x)$ in the same original metric.

Every subsequent bound is now evaluated on this identical path, including
the centered Hilbert--Schmidt, angle and signed finite-residual controls.
Retain the original monomial row $v(y)=(1,S,\ldots,S^{2q})$,
$S=k/2+iy$, in the same $H$, and the original interpolation measure
$d\mu_x(y)=((1-x)r_{1/4}^{*k}(y)+xw_h^{*k}(y))\,dy$.
For compact notation put
$n=2q+1$, $a_q=2q-1$, and define
\[
\begin{split}
 T_q(x)&=\tfrac12d_x\|U_x-W_x\|_{1,M_x},\\
 S_q(x)&=\sqrt{\left(\operatorname{Tr}(C_x^2)
       -\frac{(\operatorname{Tr}C_x)^2}{n}\right)
                    \operatorname{Tr}((U_x-W_x)^2)},\\
 z_x&=a_q\sqrt{1-e^{-F(x)/a_q}},\qquad A_q(x)=z_xd_x,\\
 J_q&=\int_0^1\min\{S_{\rm AW}(x),T_q(x),S_{\rm SP}(x),
                                      S_q(x),A_q(x)\}\,dx,\\
 I_q&=\int_0^1\min\{d_x,S_{\rm AW}(x)/z_x,T_q(x)/z_x,
                         S_{\rm SP}(x)/z_x,S_q(x)/z_x\}\,dx.
\end{split}
\tag{ACM11a}
\]
These symbols are exactly $E_q,T_q,O_q,S_q,A_q,J_q,I_q$ of
(RMT9)--(RMT13), with $E_q=S_{\rm AW}$ and $O_q=S_{\rm SP}$.
Equations (ACM1)--(ACM5) prove the operator and measure identities
needed for that substitution. In particular they retain the actual
inverse metric in the signed density $v(U_x-W_x)M_x^{-1}v^*$.

Here is the proof of the additional terms at this earlier use site.
Trace zero in (ACM10) allows subtraction of $\operatorname{Tr}C_x/n$
from $C_x$; Hilbert--Schmidt Cauchy--Schwarz gives $|F'|\leq S_q$.
The midpoint argument already proved gives $|F'|\leq T_q\leq S_{\rm SP}$.
Retain both complete $q$-entry angle lists of the minimum-section
pairs $(q-1,2q)$ and $(q,2q-1)$. The second pair has a designated
zero angle by the original ambient dimension. Each additional zero
remains in its list. The exact projection-difference spectra give
$|F'|\leq d_x\sum\sin\theta$ over the two lists. Their cosine-square
determinants give $F=\sum-\log\cos^2\theta$; applying concavity of
$t\mapsto\sqrt{1-e^{-t}}$ to the remaining $2q-1$ entries gives
$|F'|\leq A_q$, exactly as proved with every multiplicity in
(RMT11). The designated zero still contributes its displayed zero
to both scalar sums.

The denominator in (ACM11a) is strictly positive on this actual path.
Indeed equality $F=0$ makes both crossed minimum sections equal:
their Gram difference is the squared norm of their difference,
and every relative eigenvalue is at least one. Equality of their
determinants makes every such eigenvalue one. In particular the
constant remainder $1$ would be orthogonal to $\chi\mathcal P_q$.
The polynomial $\chi^{\#_k}(S)=\overline{\chi(k-\overline S)}$
has degree $q$ and equals $\overline{\chi(S)}$ on $S=k/2+iy$.
That orthogonality would imply $\int|\chi|^2d\mu_x=0$, contradicting
the positive Gram of the nonzero polynomial $\chi$. This proves
$F(x)>0$ without inserting a new hypothesis. Compactness gives a
positive minimum. All remaining functions are continuous except
the intersection dimension, whose rank strata are Borel sets
defined by matrix minors. Thus every displayed integral is finite.

The connection to the retained $(J_*,I_*)$ is an exact identity, not
an additional choice of estimate. Put $\Pi_N=P_N-Q_N$, the orthogonal
projection onto the original minimum-section space at cutoff $N$.
Substitution in (ACM5), including $Q_{q-1}=0$, gives
\[
 U_x-W_x=(\Pi_{q-1}-\Pi_{2q})+(\Pi_q-\Pi_{2q-1}).
\]
For each pair of rank-$q$ projections in $H$ of dimension $2q+1$,
the complete difference spectrum consists of the positive and negative
sine of each of its $q$ principal angles and one additional ambient
zero. All zero angles still contribute both of their zero entries.
The triangle inequality for the trace norm
therefore gives
\[
 \tfrac12\|U_x-W_x\|_{1,M_x}
 \leq\sum_{j=1}^q\sin\theta_j+\sum_{j=1}^q\sin\varphi_j
 \leq a_q\sqrt{1-e^{-F(x)/a_q}}=z_x.
 \tag{ACM11b}
\]
Here $(\theta_j)$ and $(\varphi_j)$ are respectively the two complete
angle lists just used; $\varphi_1=0$ is the designated forced zero.
The last inequality is the same concavity calculation, and no other
zero entry is removed. Consequently $T_q\leq A_q$ and the already
proved $T_q\leq S_{\rm SP}$ gives, pointwise and then by integration,
\[
\begin{split}
 J_q&=\int_0^1\min\{S_{\rm AW},T_q,S_q\}\,dx=J_*,\\
 I_q&=\int_0^1\frac{\min\{S_{\rm AW},T_q,S_q\}}{z_x}\,dx=I_*.
\end{split}
\tag{ACM11c}
\]
These are precisely the three terms in (AT21a)--(AT21c), through the
already proved identical coefficients, measures and operators.
The exact path dictionary is $B(x)=F(x)$, $a=a_q$ and $f_a(B(x))=z_x$;
it preserves the original initial value and every measure factor.
For $I_q$, the omitted entry $d_x$ is at least $T_q/z_x$ by
(ACM11b); the omitted overlap term is also at least $T_q/z_x$.
Thus both equalities retain the same original nonlinear denominator.

Both bounds are finite and measurable on $[0,1]$. Indeed the positive
Grams and their inverses vary smoothly, the eigenvalues vary
continuously, and the dimension $s_x$ is a Borel function: in any
fixed coefficient frame it is determined by the ranks of matrices
of continuous functions, whose rank strata are given by vanishing
and nonvanishing minors. Thus the triangle inequality for integration
gives the combined, original-source control
\[
 \boxed{
 |\Delta_{h,k}|\le J_*
 \le\underbrace{\int_0^1\min\{S_{\rm AW}(x),
       \tfrac12d_x\|U_x-W_x\|_{1,M_x}\}\,dx}_{\mathfrak C_*}
 \le\underbrace{\int_0^1\min\{S_{\rm AW}(x),S_{\rm SP}(x)\}\,dx}
                _{\mathfrak C_*^{\rm ov}}.}
 \tag{ACM11}
\]
The first bound preserves the exact trace norm already present in
(SP13); its proof above gives $d_x\|U_x-W_x\|_1/2\le S_{\rm SP}(x)$.
The identity maps (ACM1)--(ACM5) identify these quantities with
(AT18a), including the source density and its mass. If $b_1\le\cdots\le b_{2q+1}$ are the positive eigenvalues of
$D=M_0^{-1}M_1$ in $M_0$, the original interpolation gives
\[
 c_j(x)=\frac{b_j-1}{1-x+xb_j},\qquad
 \int_0^1c_j(x)\,dx=\log b_j.
 \tag{ACM12}
\]
The order is preserved since differentiation of the displayed
fraction with respect to $b_j$ gives $(1-x+xb_j)^{-2}>0$.
Integration of (ACM9) consequently yields
\[
 \begin{aligned}
 |\Delta_{h,k}|&\le J_*\le\mathfrak C_*
 \le\mathfrak C_*^{\rm ov}
 \le\min\left\{\mathcal L_q(D),\int_0^1S_{\rm SP}(x)\,dx\right\},\\
 \mathcal L_q(D)&=\log\frac{b_{q+2}}{b_q}
       +2\sum_{j=1}^{q-1}\log\frac{b_{2q+2-j}}{b_j}.
 \end{aligned}
 \tag{ACM13}
\]

The same pointwise minimum also retains the individually integrated
trace, centered variance and angle bounds. Termwise integration in
(ACM12), with the original order and all multiplicities, proves
\[
\begin{split}
 J_q&\leq\min\left\{\mathcal L_q(D),\int_0^1T_q(x)\,dx,
 \int_0^1S_{\rm SP}(x)\,dx,\int_0^1S_q(x)\,dx,
                                      \int_0^1A_q(x)\,dx\right\},\\
 \mathcal L_q(D)&\leq(2q-1)\log(b_{2q+1}/b_1).
\end{split}
\tag{ACM13a}
\]
The last inequality bounds each of the $1+2(q-1)$ logarithmic factors
in the unchanged (ACM13) by $\log(b_{2q+1}/b_1)$. It is therefore
valid also for $q=1$, where the original sum is empty.

The empty sum at $q=1$ is zero. No universal ordering of the two
refinements is used: each bounds the same signed integrand, hence
their pointwise minimum does. Their full original arithmetic
moments, relative eigenvalues and overlap remain the actual terms
to be evaluated along the growing packet family.

The actual moment interpolation above also satisfies the complete
proof (AT21a)--(AT21e). With $a=2q-1$ its recursive additive
minimum also retains the independently proved centered
HS term $S_{\rm HS}$ in (AT21a), so $J_*\le\mathfrak C_*$.
The exact trace
norm is at most both the overlap and full angle bounds. The
nonlinear radius
$I_*=\int_0^1\min(S_{\rm spec},S_{\rm tr},S_{\rm HS})/
[a\sqrt{1-e^{-B(x)/a}}]\,dx$
is defined on the same strictly positive four-volume path, with
$I_*\le\log\kappa$. Thus the current endpoint consequence here
is the simultaneous interval $\mathfrak b_*^-\le B(1)\le
\mathfrak b_*^+$ of (AT21e), not only its additive corollary.
The identity of the coefficient maps and original measures proved
above verifies its hypotheses; no purity or tensor-uniform bound
is supplied by this finite transport.

The following signed construction further intersects this same retained
additive/nonlinear interval. By (ACM11c) its $(J_q,I_q)$ are literally
$(J_*,I_*)$; hence no earlier endpoint or overlap bound is discarded.

The nonlinear coordinate also receives every one of these controls.
For $B>0$ put
\[
 \mathscr H_{a_q}(B)=2\operatorname{arcosh}(e^{B/(2a_q)}),\quad
 \Psi_{a_q}(t)=2a_q\log\cosh(t/2)\quad(t\geq0).
 \tag{ACM14}
\]
Their two inverse composites on the positive half-lines are identities.
Direct differentiation gives
\[
 \mathscr H_{a_q}'(B)=\frac{1}{a_q\sqrt{1-e^{-B/a_q}}}.
\]
Dividing the five proved derivative bounds
by $z_x$ and integrating therefore gives
$|\mathscr H_{a_q}(F(1))-\mathscr H_{a_q}(F(0))|\leq I_q$.
Moreover $I_q\leq\int d_x=\log(b_n/b_1)$.

For every integer $p\geq0$ retain the following specified finite
signed construction, on the same source operators:
\[
\begin{split}
 B_x^\circ&=(1-x)I+xD,\qquad
 \alpha_x=\tfrac12(2(1-x)+x(b_1+b_n)),\qquad
 H_x^\circ=I-B_x^\circ/\alpha_x,\\
 C_x^{[p]}&=\alpha_x^{-1}\sum_{r=0}^p(H_x^\circ)^r(D-I),\\
 R_x^{[p]}&=(H_x^\circ)^{p+1}C_x,
 \qquad E_x^{[p]}=R_x^{[p]}-\operatorname{Tr}(R_x^{[p]})I/n,\\
 j_p&=\int_0^1\operatorname{Tr}(C_x^{[p]}(U_x-W_x))\,dx,\\
 e_p&=\int_0^1\sqrt{\operatorname{Tr}((E_x^{[p]})^2)
                         \operatorname{Tr}((U_x-W_x)^2)}\,dx.
\end{split}
\tag{ACM15}
\]
The finite geometric identity proves $C_x-C_x^{[p]}=R_x^{[p]}$.
Every factor is a real function of $D$, hence self-adjoint in $M_x$.
Trace zero removes the centered scalar, so Hilbert--Schmidt
Cauchy--Schwarz proves $j_p-e_p\leq\Delta_{h,k}\leq j_p+e_p$.
Writing $\vartheta=(b_n-b_1)/(b_n+b_1)<1$ and
$K_D=\sum_{j=1}^n(|b_j-1|/\min\{1,b_j\})^2$, the full residual
eigenvalues give
$\operatorname{Tr}((E_x^{[p]})^2)\leq\vartheta^{2p+2}K_D$.
Both $U_x,W_x$ dominate their range projections, whose intersection
has dimension at least one, so $\operatorname{Tr}(U_xW_x)\geq1$.
Equation (ACM10) then gives
$e_p\leq\vartheta^{p+1}\sqrt{K_D(8q-6)}\to0$.
This retains the exact signed center and its controlled residual,
with no assumption of uniformity in the tensor order.

Put $F_0=F(0)$ and $H_0=\mathscr H_{a_q}(F_0)$. For every $p\geq0$
define the actual simultaneous endpoint functions
\[
\begin{split}
 \beta_p^-&=\max\{0,F_0-J_q,
       \Psi_{a_q}(\max\{0,H_0-I_q\}),F_0+j_p-e_p\},\\
 \beta_p^+&=\min\{F_0+J_q,
       \Psi_{a_q}(H_0+I_q),F_0+j_p+e_p\}.
\end{split}
\tag{ACM16}
\]
The preceding additive, nonlinear and signed inequalities prove
$\beta_p^-\leq F(1)\leq\beta_p^+$ by taking their intersection.
Both original angle lists, source masses, relation maps and arithmetic
Taylor unit remain fixed. At $D=cI$ all diameter and residual terms
vanish: each quotient volume is multiplied by the same scalar power
$q$, which cancels with the four signs. The formulas then give
$F(1)=F_0$ exactly. Numerical integration of a signed center must
retain its own integration error; (ACM15) specifies exact integrals.

\subsection{The literal graph augmentation on this finite coefficient source}

The following additional calculation uses the full quartet domain of
PGB.1: $\gamma>2$, $0<\delta<1/2$, $m,k\geq1$.
It leaves the preceding Gamma-to-arithmetic formulas on their entire
stated domain. Here $\chi$ is the literal monic polynomial
\[
 \chi(S)=\prod_{a,b=0}^{k}
 \bigl(S-k/2-(2a-k)\delta-i(2b-k)\gamma\bigr)^{1+k(m-1)}.
 \tag{ACM17}
\]
Thus (HT2) and PGB.1 have identical roots, multiplicities, leading
coefficient, increasing remainder frame and remainder map; their
coefficient identification is the identity, including every nilpotent
primary coordinate. Retain $g=2\xi$, $v_h=g/h$, $\Phi'=h$,
$\Phi(0)=0$, $s_i=1/2+it_i$, $\Sigma=\sum_i s_i$ and
$\tau_k=\sum_i\Phi(s_i)$. The integration measure in the full
source is $\prod_i dt_i/(2\pi)$, and
\[
 \begin{aligned}
 (\mathcal VP)(\mathbf t)&=\Bigl(\prod_i v_h(s_i)\Bigr)P(\Sigma),\\
 L_{\rm obs}x&=\Bigl(\prod_i v_h(s_i)\Bigr)
  \sum_{b=0}^{q-1}x_b\operatorname{rem}_{h(s_1),\ldots,h(s_k)}
                                      (\tau_k\Sigma^b),\\
 D_{\rm def}P&=\tau_k\mathcal VP-L_{\rm obs}J_{2q}P,
 \qquad H=\mathcal P_{2q},\qquad M^{\rm ar}=\mathcal V^*\mathcal V.
 \end{aligned}
 \tag{ACM18}
\]
The symbols $L_{\rm obs}$ and $D_{\rm def}$ denote exactly $L$ and
$D_{2q}$ of PGB.2; neither is the circle circumference or the
positive relative matrix below. PGB.2 and its leading-coefficient
argument prove that these columns lie in the stated Hilbert source
and that $D_{\rm def}$ is injective. In particular every following
finite Gram exists and is positive. On the unchanged coefficient
space define the actual affine path
\[
 M_x^{\rm gr}=M^{\rm ar}+xD_{\rm def}^*D_{\rm def},\quad
 0\leq x\leq1,\qquad
 \mathcal V_x^{\rm gr}P=(\mathcal VP,\sqrt{x}D_{\rm def}P).
 \tag{ACM19}
\]
Its Gram is exactly $M_x^{\rm gr}$, by the orthogonal direct-sum
norm; no mixed term was omitted. At $x=1$ this is the unchanged
graph of PGB.2. The restriction to $\mathcal P_N$ uses the literal
$I_N$ and has defect $D_{\rm def}I_N=D_N$: monic remainders
satisfy $J_{2q}I_N=J_N$, proving this equality directly in (ACM18).
Its relation columns are $I_NB_N^{\rm AW}$, exactly as in (ACM3).

Use (ACM4)--(ACM5) with this new, explicitly specified form and put
$U_x^{\rm gr}=Q_{2q-1}+Q_{2q}-Q_q$ and
$W_x^{\rm gr}=P_{2q-1}-P_{q-1}+P_{2q}-P_q$.
Let $G_N(x)$ denote the quotient Gram of its restriction and set
\[
 \begin{aligned}
 F_{\rm gr}(x)&=\log\det G_{q-1}(x)+\log\det G_q(x)
              -\log\det G_{2q-1}(x)-\log\det G_{2q}(x),\\
 C_x^{\rm gr}&=(M_x^{\rm gr})^{-1}D_{\rm def}^*D_{\rm def},
 \qquad A_x^{\rm gr}=U_x^{\rm gr}-W_x^{\rm gr},\\
 \Gamma_k^{\rm gr}&=F_{\rm gr}(1)-F_{\rm gr}(0)
                 =\widehat{\mathcal B}_k-\mathcal B_k.
 \end{aligned}
 \tag{ACM20}
\]
Here $\widehat{\mathcal B}_k$ is the four-endpoint logarithm of
$\widehat G_N$ in PGB.26, and $\mathcal B_k$ is that of the
original arithmetic source. To prove the derivative at this use
site, the minimum section $R_N(x)$ satisfies $J_NR_N=I$ and
$B_N^*M_N(x)R_N=0$. Differentiating the first equality puts
$\dot R_N$ in the relation space. The two cross terms in
$d(R_N^*M_NR_N)/dx$ therefore vanish by the second equality,
and $\dot G_N=R_N^*\dot M_NR_N$. Jacobi's determinant formula
and finite trace cyclicity give
$d\log\det G_N/dx=\operatorname{Tr}((P_N-Q_N)C_x^{\rm gr})$.
Taking the four signs proves
\[
 F_{\rm gr}'(x)=\operatorname{Tr}(A_x^{\rm gr}C_x^{\rm gr}).
 \tag{ACM21}
\]
The flag inclusions imply that each difference of nested
orthogonal projections is an orthogonal projection. In particular
$Q_q$ has rank one and is contained in both $Q_{2q-1}$ and
$Q_{2q}$. The original dimensions give the complete spectra
$0^{[q]},1^{[2]},2^{[q-1]}$ for $U_x^{\rm gr},W_x^{\rm gr}$.
Their traces agree, so $\operatorname{Tr}A_x^{\rm gr}=0$ and
\[
 \operatorname{Tr}((A_x^{\rm gr})^2)
       =8q-4-2\operatorname{Tr}(U_x^{\rm gr}W_x^{\rm gr}).
 \tag{ACM22}
\]
These arguments use positive finite forms and the same flags;
they do not use a scalar moment representation of the graph norm.
Hence the proofs at (ACM9)--(ACM11a) apply to their spectral,
trace-norm, overlap and centered Hilbert--Schmidt terms. Explicitly
let $c_1^{\rm gr}\leq\cdots\leq c_n^{\rm gr}$ be the full
eigenvalue list of $C_x^{\rm gr}$, $n=2q+1$, and let
$d_x^{\rm gr}=c_n^{\rm gr}-c_1^{\rm gr}$ and
$s_x^{\rm gr}=\dim(\operatorname{ran}U_x^{\rm gr}
\cap\operatorname{ran}W_x^{\rm gr})$. Set
\[
 \begin{aligned}
 E_{\rm gr}(x)&=c_{q+2}^{\rm gr}-c_q^{\rm gr}
  +2\sum_{j=1}^{q-1}(c_{2q+2-j}^{\rm gr}-c_j^{\rm gr}),\\
 T_{\rm gr}(x)&=\tfrac12 d_x^{\rm gr}
                   \|A_x^{\rm gr}\|_{1,M_x^{\rm gr}},\\
 O_{\rm gr}(x)&=\tfrac12d_x^{\rm gr}\min\left\{
 4q-2s_x^{\rm gr},\sqrt{n\operatorname{Tr}((A_x^{\rm gr})^2)}
                                                   \right\},\\
 S_{\rm gr}(x)&=\sqrt{\left(\operatorname{Tr}((C_x^{\rm gr})^2)
 -\frac{(\operatorname{Tr}C_x^{\rm gr})^2}{n}\right)
                   \operatorname{Tr}((A_x^{\rm gr})^2)},\\
 J_{\rm gr}&=\int_0^1\min\{E_{\rm gr},T_{\rm gr},O_{\rm gr},
                                                S_{\rm gr}\}\,dx.
 \end{aligned}
 \tag{ACM23}
\]
Rank-stratum measurability and continuity of the other entries
prove integrability exactly as before. The four pointwise trace
proofs give $|\Gamma_k^{\rm gr}|\leq J_{\rm gr}$.

For a specified integer $p\geq0$, this same path also has an
actual signed finite approximation. Put
\[
 \begin{aligned}
 D^{\rm gr}&=(M^{\rm ar})^{-1}
               (M^{\rm ar}+D_{\rm def}^*D_{\rm def}),
 \qquad 1<b_1^{\rm gr}\leq\cdots\leq b_n^{\rm gr},\\
 B_x^{\rm gr}&=(1-x)I+xD^{\rm gr},\qquad
 \alpha_x^{\rm gr}=1-x+\tfrac{x}{2}(b_1^{\rm gr}+b_n^{\rm gr}),\\
 H_x^{\rm gr}&=I-B_x^{\rm gr}/\alpha_x^{\rm gr},\qquad
 C_x^{[p],\rm gr}=(\alpha_x^{\rm gr})^{-1}
     \sum_{r=0}^{p}(H_x^{\rm gr})^r(D^{\rm gr}-I),\\
 R_x^{[p],\rm gr}&=(H_x^{\rm gr})^{p+1}C_x^{\rm gr},\qquad
 E_x^{[p],\rm gr}=R_x^{[p],\rm gr}
                 -\operatorname{Tr}(R_x^{[p],\rm gr})I/n,\\
 j_{p,\rm gr}&=\int_0^1\operatorname{Tr}
                         (C_x^{[p],\rm gr}A_x^{\rm gr})\,dx,\\
 e_{p,\rm gr}&=\int_0^1\sqrt{
 \operatorname{Tr}((E_x^{[p],\rm gr})^2)
 \operatorname{Tr}((A_x^{\rm gr})^2)}\,dx.
 \end{aligned}
 \tag{ACM24}
\]
The strict lower bound for $b_1^{\rm gr}$ follows from injectivity
of $D_{\rm def}$ in this finite dimension. The finite geometric
identity gives $C_x^{\rm gr}-C_x^{[p],\rm gr}=R_x^{[p],\rm gr}$.
All these factors are real functions of the same $D^{\rm gr}$,
so they are self-adjoint for $M_x^{\rm gr}$. Pairing the last
identity with the trace-zero $A_x^{\rm gr}$ and applying
Hilbert--Schmidt Cauchy--Schwarz proves
\[
 j_{p,\rm gr}-e_{p,\rm gr}\leq\Gamma_k^{\rm gr}
                     \leq j_{p,\rm gr}+e_{p,\rm gr}.
 \tag{ACM25}
\]
For completeness set
\[
 \vartheta_{\rm gr}=
 \frac{b_n^{\rm gr}-b_1^{\rm gr}}{b_n^{\rm gr}+b_1^{\rm gr}},
 \qquad
 K_{\rm gr}=\sum_{j=1}^{n}
 \left(\frac{|b_j^{\rm gr}-1|}{\min\{1,b_j^{\rm gr}\}}\right)^2.
\]
The eigenvalue formula
$c_j^{\rm gr}(x)=(b_j^{\rm gr}-1)/(1-x+xb_j^{\rm gr})$
shows $\|H_x^{\rm gr}\|\leq\vartheta_{\rm gr}<1$ and
$\operatorname{Tr}((E_x^{[p],\rm gr})^2)
\leq\vartheta_{\rm gr}^{2p+2}K_{\rm gr}$.
The range intersection has dimension at least one. Since each
of $U_x^{\rm gr},W_x^{\rm gr}$ dominates its range projection,
the positive-product trace inequality gives
$\operatorname{Tr}(U_x^{\rm gr}W_x^{\rm gr})\geq1$.
Together with (ACM22), this proves
\[
 e_{p,\rm gr}\leq\vartheta_{\rm gr}^{p+1}
                   \sqrt{K_{\rm gr}(8q-6)}\longrightarrow0.
 \tag{ACM26}
\]
Thus the two actual signed endpoints
\[
 g_{p,\rm gr}^-:=\max\{-J_{\rm gr},j_{p,\rm gr}-e_{p,\rm gr}\},
 \qquad
 g_{p,\rm gr}^+:=\min\{J_{\rm gr},j_{p,\rm gr}+e_{p,\rm gr}\}
 \tag{ACM27}
\]
satisfy $g_{p,\rm gr}^-\leq\Gamma_k^{\rm gr}
\leq g_{p,\rm gr}^+$. This calculation uses no division by
the angle coordinate of the new graph form. The moment-path
nonlinear formulas (ACM14)--(ACM16) retain their original scope.
All integrals in (ACM23)--(ACM27) are exact integrals of the
displayed finite matrices; numerical evaluation must add its
own certified integration error.

\subsection{The complete original relation rows in the graph endpoint control}

Keep exactly the domain and coefficient identity of (ACM17)--(ACM20).
The symbols in this paragraph include a superscript $\mathrm{rel}$
so that the monic relation polynomial is never confused with a
consecutive quotient-volume ratio. They are the literal objects of
(PGRT.2)--(PGRT.8):
\[
\begin{gathered}
 r_k(u)=\int_{\sum t_i=u}(1+|\tau_k|^2)\prod_iw_h(t_i),\qquad
 F_b=\operatorname{rem}_{h(s_1),\ldots,h(s_k)}(\tau_k\Sigma^b),\\
 z_b(u)=\int_{\sum t_i=u}\overline{\tau_k}F_b\prod_iw_h(t_i),
 \quad g_x=z(u)x/r_k(u),\quad
 F^0x=(sx)(k/2+iu)-g_x(u),\\
 d_j^{\rm rel}=\chi p_{j+1-q},\qquad
 a_j^{\rm rel}=\int|d_j^{\rm rel}(k/2+iu)|^2r_k(u)\,du>0,\\
 (\alpha_j^{\rm rel})_b=
 \int\overline{d_j^{\rm rel}(k/2+iu)}
       ((k/2+iu)^b r_k(u)-z_b(u))\,du,
 \quad j\geq q-1,\quad 0\leq b<q.
\end{gathered}\tag{ACM28}
\]
Here $s:E\to\mathcal P_{q-1}$ is the original monic remainder
section, $p_n$ is the unique monic polynomial of degree $n$
orthogonal to lower polynomials for $|\chi|^2r_k\,du$, and the
fibre integral eliminates $t_k$ with its original Lebesgue measure.
For $k=1$ it is evaluation. The constant factors are
$w_h(t)=|v_h(1/2+it)|^2/(2\pi)$ throughout. Each coefficient is
absolutely integrable by the proved polynomial exponential moments
and Cauchy--Schwarz applied to $F^0$. Expanding the full polynomials
in these integrals gives the finite products of the original
one-factor moments in (PGS.35) and (PGS.38), without dropping their
conjugate or cross terms. The coefficients of $p_n$ are obtained
from its positive finite lower-polynomial Gram system. This is
the exact finite arithmetic evaluation of each row.

The complete weighted density proof (PGRT.3)--(PGRT.6) gives
\[
 \widehat G_N-\Omega_k
  =\sum_{j=N}^{\infty}
       (a_j^{\rm rel})^{-1}(\alpha_j^{\rm rel})^*
                                  \alpha_j^{\rm rel},
 \qquad \Omega_k=K_{\rm low}+D_{\rm mix}^*D_{\rm mix}>0.
 \tag{ACM29}
\]
The series converges in the original finite matrix operator norm.
Specifically, the relation columns are a complete orthogonal family
in $L^2(r_kdu)$; the coefficient of $F^0x$ on its $j$th column is
$(a_j^{\rm rel})^{-1}\alpha_j^{\rm rel}x$. Removing the first
$N-q+1$ columns is precisely minimization over the original
relations $\chi\mathcal P_{N-q}$. Parseval gives the displayed
quadratic forms, and polarization gives their matrices. The trace
of the remaining positive tail is the sum of the squared tails
of the $q$ fixed frame vectors, hence tends to zero and bounds
its operator norm. The positive completed term and its mixed
realization are the same as (HC22), with full proof in PGMB.

For $q-1\leq j\leq2q-1$ put
\[
\begin{aligned}
 l_j^{\rm rel}
   &=(a_j^{\rm rel})^{-1}\alpha_j^{\rm rel}
             \widehat G_{q-1}^{-1}(\alpha_j^{\rm rel})^*,\\
 u_j^{\rm rel}
   &=(a_j^{\rm rel})^{-1}\alpha_j^{\rm rel}
             \widehat G_{2q}^{-1}(\alpha_j^{\rm rel})^*,\\
 Z_{2q}^{\rm mix}x&=(D_{\rm mix}x,f_{2q}x),\\
 u_j^{\rm rel}
   &=(a_j^{\rm rel})^{-1}\alpha_j^{\rm rel}
                     K_{\rm low}^{-1}(\alpha_j^{\rm rel})^*\\
   &\quad-(a_j^{\rm rel})^{-1}\left\|
 (I+Z_{2q}^{\rm mix}K_{\rm low}^{-1}(Z_{2q}^{\rm mix})^*)^{-1/2}
 Z_{2q}^{\rm mix}K_{\rm low}^{-1}(\alpha_j^{\rm rel})^*
                                      \right\|^2.
\end{aligned}\tag{ACM30}
\]
The domain of $Z_{2q}^{\rm mix}$ is the original $E$, and its
ordered target is the mixed weighted Hilbert space of (HC22)
followed by $L^2(r_kdu)$. In particular it retains the finite
tail as well as the mixed component. To prove the subtraction,
write $K=K_{\rm low}$ and $Z=Z_{2q}^{\rm mix}$. The identity
$\widehat G_{2q}=K+Z^*Z$ gives
\[
 (K+Z^*Z)^{-1}
 =K^{-1}-K^{-1}Z^*(I+ZK^{-1}Z^*)^{-1}ZK^{-1}.
\]
Multiplication by $K+Z^*Z$ cancels the leftover factor
$Z^*ZK^{-1}-Z^*(I+ZK^{-1}Z^*)(I+ZK^{-1}Z^*)^{-1}ZK^{-1}$.
The target inverse exists because it is the identity off the
finite-dimensional range of $Z$ and is positive with eigenvalues
at least one on that range. Pairing with the specified row proves
(ACM30). If $k\geq2$ and $\alpha_j^{\rm rel}\ne0$, the
subtracted squared norm is strictly positive: $D_{\rm mix}$ is
injective by the full fixed-sum calculation PGM.6--9, and the
displayed inverse square root has zero kernel. For $k=1$ its
mixed component is zero but its finite tail is retained. A zero
row gives zero in both terms. These statements retain the original
strictness domain and the whole negative term.

Use the endpoint weights $\omega_{q-1}=\omega_{2q-1}=1$ and
$\omega_j=2$ for $q\leq j\leq2q-2$, and define
\[
\begin{gathered}
 L_k^{\rm rel}=\sum_{j=q-1}^{2q-1}\omega_j\log(1+l_j^{\rm rel}),
 \qquad
 U_{k,2q}^{\rm rel}=\sum_{j=q-1}^{2q-1}\omega_j\log(1+u_j^{\rm rel}),\\
 L_k^{\rm rel}\leq\widehat{\mathcal B}_k
                                  \leq U_{k,2q}^{\rm rel},\\
 L_k^{\rm rel}-g_{s,\rm gr}^{+}\leq\mathcal B_k
                       \leq U_{k,2q}^{\rm rel}-g_{s,\rm gr}^{-}
 \qquad(s\geq0).
\end{gathered}\tag{ACM31}
\]
To verify every index, (ACM29) gives
$\widehat G_j=\widehat G_{j+1}
 +(a_j^{\rm rel})^{-1}(\alpha_j^{\rm rel})^*\alpha_j^{\rm rel}$.
The rank-one determinant identity therefore gives the $j$th
decrement $\log(1+(a_j^{\rm rel})^{-1}\alpha_j^{\rm rel}
\widehat G_{j+1}^{-1}(\alpha_j^{\rm rel})^*)$.
Since $\widehat G_{2q}\preceq\widehat G_{j+1}
\preceq\widehat G_{q-1}$, inverse order bounds its argument
between $1+l_j^{\rm rel}$ and $1+u_j^{\rm rel}$.
Inverse order follows by congruence with the smaller positive
form, inversion of its positive eigenvalues, and reversal of
that congruence. Summing the decrements from $q-1$ to $2q-2$
and from $q$ to $2q-1$ gives precisely the displayed weights.
The full quartet in (ACM17) has $q\geq4$, so the listed endpoints
and interior have these literal ranges. Finally (ACM20) says
$\mathcal B_k=\widehat{\mathcal B}_k-\Gamma_k^{\rm gr}$.
Subtracting (ACM27) yields the last line, with its ordered negative
$g_{s,\rm gr}^{-}$ term intact. This is a finite relation-row
bound on the same arithmetic scalar; it does not replace the
exact value of $\widehat{\mathcal B}_k$ when that value is retained.


\paragraph{All retained relation correlations in the same inverse.}
The complete calculation PGRC.1--18 now sharpens the row enclosure
(ACM31) by evaluating its full cross terms. Keep
$H=\widehat G_{2q}$, $I=\{q-1,\ldots,2q-1\}$ and put
$C_{ij}=\alpha_i^{\rm rel}H^{-1}(\alpha_j^{\rm rel})^*$.
For an ordered subset $J\subset I$ let $A_J$ have the corresponding
original rows, let $D_J=\operatorname{diag}(a_i^{\rm rel}:i\in J)$,
and define
\[
\begin{gathered}
 H_J=H+A_J^*D_J^{-1}A_J,\quad B_J=D_J+C_{JJ},\\
 V_J:(E,H_J)\longrightarrow(E,H)\oplus(\mathbb C^J,D_J^{-1}),
 \quad V_Jx=(x,A_Jx),\\
 V_J^\dagger(y,z)=H_J^{-1}(Hy+A_J^*D_J^{-1}z),\\
 V_J^\dagger V_J=I_E,\quad\ker V_J=0,\quad
 \ker V_J^\dagger=\{(y,z):Hy+A_J^*D_J^{-1}z=0\}.
\end{gathered}\tag{ACM32}
\]
The norm expansion proves the Gram $H_J$ and the adjoint;
multiplication proves the inverse on the image and both kernel
descriptions. Empty subsets mean the zero-dimensional coordinate
space and its unique empty inverses. For every subset, direct
multiplication also proves
$H_J^{-1}=H^{-1}-H^{-1}A_J^*B_J^{-1}A_JH^{-1}$:
the extra middle factor is
$D_J^{-1}-B_J^{-1}-D_J^{-1}C_{JJ}B_J^{-1}=0$
because $I+D_J^{-1}C_{JJ}=D_J^{-1}B_J$.

Let $J_j=\{j+1,\ldots,2q-1\}$. For $J\subset J_j$ put
$r_j(J)=C_{jJ}B_J^{-1}C_{Jj}$ and
$v_j(J)=(C_{jj}-r_j(J))/a_j^{\rm rel}$.
The inverse identity proves
$v_j(J)=\alpha_j^{\rm rel}H_J^{-1}(\alpha_j^{\rm rel})^*/a_j^{\rm rel}
\geq0$. Telescoping (ACM29) identifies
$H_{J_j}=\widehat G_{j+1}$. Thus the exact decrement parameter is
\[
 \eta_j=
 \frac{C_{jj}-C_{jJ_j}(D_{J_j}+C_{J_jJ_j})^{-1}C_{J_jj}}
      {a_j^{\rm rel}}.
 \tag{ACM33}
\]
For disjoint $J,K\subset J_j$, keep the remaining block and row
\[
\begin{aligned}
 S_{K\mid J}&=D_K+C_{KK}-C_{KJ}B_J^{-1}C_{JK}
             =D_K+A_KH_J^{-1}A_K^*\succeq D_K>0,\\
 e_{j,K\mid J}&=C_{jK}-C_{jJ}B_J^{-1}C_{JK}
                =\alpha_j^{\rm rel}H_J^{-1}A_K^*,\\
 r_j(J\cup K)-r_j(J)&=e_{j,K\mid J}S_{K\mid J}^{-1}e_{j,K\mid J}^*,\\
 0\leq r_j(J\cup K)-r_j(J)
 &\leq\sum_{i\in K}|(e_{j,K\mid J})_i|^2/a_i^{\rm rel}.
\end{aligned}\tag{ACM34}
\]
Substitution of the proved inverse gives the first two lines.
For interleaved disjoint subsets, first use the exact coordinate map
$P_{J,K}:\mathbb C^{J\cup K}\to\mathbb C^J\oplus\mathbb C^K$,
$z\mapsto(z|_J,z|_K)$, with inverse reassembling the coordinates
in their inherited order. This is an isometry for the retained
diagonal lengths: its two squared norms sum to the original
sum over $J\cup K$. On coefficient matrices it is a permutation,
so $P_{J,K}(D_{J\cup K}+C_{J\cup K,J\cup K})P_{J,K}^{-1}$
is exactly the displayed two-block matrix. Conjugating the inverse
and applying the same restriction to the row leaves the scalar
$r_j(J\cup K)$ unchanged. To prove the third line, solve the
resulting two block equations
$B_Jx+C_{JK}y=C_{Jj}$ and
$C_{KJ}x+(D_K+C_{KK})y=C_{Kj}$.
The first gives $x=B_J^{-1}(C_{Jj}-C_{JK}y)$ and the second
then gives $y=S_{K\mid J}^{-1}e_{j,K\mid J}^*$.
Pairing with $(C_{jJ},C_{jK})$ proves exactly the displayed
quadratic correction. Inverse order from $S_{K\mid J}\succeq D_K$
gives the final inequality. In particular all cross subtractions
are retained even when they vanish.

Choose actual finite subsets $J(j)\subset J_j$, let
$K(j)=J_j\setminus J(j)$, and set
\[
\begin{gathered}
 \overline\eta_j=v_j(J(j)),\qquad
 \underline\eta_j=\max\left\{l_j^{\rm rel},\,
 v_j(J(j))-\sum_{i\in K(j)}
  \frac{|(e_{j,K(j)\mid J(j)})_i|^2}{a_j^{\rm rel}a_i^{\rm rel}}
                                                   \right\},\\
 L_k^{\rm corr}(J)=\sum_{j=q-1}^{2q-1}\omega_j\log(1+\underline\eta_j),
 \quad
 U_k^{\rm corr}(J)=\sum_{j=q-1}^{2q-1}\omega_j\log(1+\overline\eta_j),\\
 L_k^{\rm rel}\leq L_k^{\rm corr}(J)
 \leq\widehat{\mathcal B}_k\leq U_k^{\rm corr}(J)
 \leq U_{k,2q}^{\rm rel}.
\end{gathered}\tag{ACM35}
\]
The last line of (ACM34) with $J\cup K=J_j$ proves
$0\leq\underline\eta_j\leq\eta_j\leq\overline\eta_j$.
Moreover $\underline\eta_j\geq l_j^{\rm rel}$ and
$\overline\eta_j\leq C_{jj}/a_j^{\rm rel}=u_j^{\rm rel}$.
Increasing logarithms and the unchanged endpoint weights prove
all the inequalities in (ACM35). If every $J(j)=J_j$, the remaining
block is empty and both new bounds equal the exact graph volume.
Insertion of one further index $b$ subtracts exactly
\[
 \frac{|C_{jb}-C_{jJ}B_J^{-1}C_{Jb}|^2}
 {a_j^{\rm rel}(a_b^{\rm rel}+C_{bb}-C_{bJ}B_J^{-1}C_{Jb})}
\]
from $v_j(J)$, by (ACM34). Its denominator is positive. Thus
successive upper bounds decrease and become exact after at most
$q$ later-row insertions for each row; retaining the maximum
of all previous valid lower bounds makes that retained sequence
increasing as well.

In particular the exact $a_j^{\rm rel}\eta_j$ is the nonnegative
$\alpha_j^{\rm rel}K_{\rm low}^{-1}(\alpha_j^{\rm rel})^*$
minus the full squared norm in (ACM30) minus
$C_{jJ_j}(D_{J_j}+C_{J_jJ_j})^{-1}C_{J_jj}$.
The latter vanishes exactly when $C_{J_jj}=0$, because its
inverse matrix is positive. This proves that both negative
corrections occur in the same original inverse quantity.
Their signs and equality cases are retained. The matrix $C$
in the correlation subtraction is computed using
$\widehat G_{2q}^{-1}$, with its full mixed term; it is never
replaced there by $K_{\rm low}^{-1}$.

At every independent coefficient face $(W,S;A)$ the coefficient
inclusions and zero entries act componentwise on these same linear source maps.
Retain the original source $F_h$ with $\mathcal M F_h=g/h$,
$D_i=-x_i\partial_{x_i}$, $h(D)F_h=\Theta\phi_*$, and
$\mathcal V P=P(D_1+\cdots+D_k)F_h^{\otimes k}$ from (ISM1).
The full Taylor unit in that observation is unchanged.
The original minimum-section difference is a relation $\chi Q$;
successive monic division constructs
$\chi(\sum s_i)Q(\sum s_i)=\sum_i h(s_i)Q_i(\mathbf s)$.
Its original primitive is
$\sum_i(-1)^{i-1}Q_i(D_1,\ldots,D_k)
 (F_h^{\otimes(i-1)}\otimes\phi_*\otimes F_h^{\otimes(k-i)})$.
The tensor differential contributes the same sign and therefore
recovers the relation. The differential operators commute with the original tensor
differential, and $h(D_i)F_h=\Theta\phi_*$, proving that its image
is exactly $\mathcal V(\chi Q)$. For the original two-leg complex,
use the degree-zero map $j_+(v)=(v,0)$ and degree-one identity into
$[V\oplus\widehat V\to\mathscr B]$, whose differential is
$\Theta(\phi-\mathcal F^{-1}\widehat\psi)$.
The tensor cochain map preserves every sign and gives the same top
boundary. The one-factor two-leg differential has the entire diagonal kernel
$(v,\mathcal Fv)$: its zero equation is equivalent to
$\phi=\mathcal F^{-1}\widehat\psi$ by injectivity of $\Theta$,
and every such diagonal pair has zero differential.
When the two compared top representatives lie in an admitted
proper subcomplex $\mathsf C_{W,S;A}^\bullet$ of the fixed full
coefficient carrier, their residual is
the class in
$(d\mathsf C_{\rm full}^{k-1}\cap\mathsf C_{W,S;A}^k)/
d\mathsf C_{W,S;A}^{k-1}$, with its complete two-leg primitive retained. The inclusion-induced
map sends a numerator representative to its class in
$\mathsf C_{W,S;A}^k/d\mathsf C_{W,S;A}^{k-1}$. Its image is exactly
the kernel of the map to the full top quotient: mapping to zero there
means being a full boundary. Its inverse uses that same representative,
and both are well-defined because their denominator is the identical
$d\mathsf C_{W,S;A}^{k-1}$. This proves the stated kernel isomorphism,
as in the complete coefficient and source maps (ISM38)--(ISM43).
Consequently none of the metric estimates deletes this residual,
the full diagonal degree-zero source, or an empty mask's outer label.

